\documentclass[12pt]{article}
\usepackage{amsmath,amssymb,amsfonts}
\usepackage[margin=1in]{geometry}

\begin{document}

\textbf{Chapter 6, Problem 37.} In Section 5.10, we defined the Hermite polynomials in terms of the generating function
\[
e^{-t^2+2xt} = \sum_{n=0}^{\infty}\frac{H_n(x)}{n!}t^n, \quad t > 0,
\]
and deduced from this the Rodrigues formula:
\[
H_n(x) = (-1)^n e^{x^2}\frac{d^n}{dx^n}e^{-x^2}.
\]

Now prove the converse: that Rodrigues' formula implies the generating function. Enroute, establish the integral representation
\[
H_n(x) = \frac{n!}{2\pi i}\oint_C \frac{2^n e^{-t^2+2xt}}{t^{n+1}}\,dt \cdot (\text{or similar}),
\]
where $C$ encloses the point $x$.

\bigskip
\textbf{Solution.}

\textbf{Step 1: Integral representation.}

From Rodrigues' formula $H_n(x) = (-1)^ne^{x^2}\frac{d^n}{dx^n}e^{-x^2}$, apply Cauchy's formula for the $n$-th derivative:
\[
\frac{d^n}{dx^n}e^{-x^2} = \frac{n!}{2\pi i}\oint_C\frac{e^{-t^2}}{(t-x)^{n+1}}\,dt,
\]
where $C$ encloses $t = x$.

Therefore:
\[
H_n(x) = (-1)^ne^{x^2}\cdot\frac{n!}{2\pi i}\oint_C\frac{e^{-t^2}}{(t-x)^{n+1}}\,dt = \frac{(-1)^nn!}{2\pi i}\oint_C\frac{e^{x^2-t^2}}{(t-x)^{n+1}}\,dt.
\]

Let $u = t - x$, so $t = u + x$ and $t^2 = u^2+2ux+x^2$:
\[
x^2 - t^2 = -u^2 - 2ux.
\]
\[
H_n(x) = \frac{(-1)^nn!}{2\pi i}\oint\frac{e^{-u^2-2ux}}{u^{n+1}}\,du.
\]

Now substitute $u = -s$:
\[
H_n(x) = \frac{(-1)^nn!}{2\pi i}\oint\frac{e^{-s^2+2sx}}{(-s)^{n+1}}(-ds) = \frac{n!}{2\pi i}\oint\frac{e^{-s^2+2sx}}{s^{n+1}}\,ds.
\]

So:
\[
\boxed{H_n(x) = \frac{n!}{2\pi i}\oint_C\frac{e^{-s^2+2sx}}{s^{n+1}}\,ds.}
\]

\textbf{Step 2: Derive the generating function.}

The generating function is:
\[
\sum_{n=0}^{\infty}\frac{H_n(x)}{n!}t^n = \sum_{n=0}^{\infty}\frac{t^n}{2\pi i}\oint\frac{e^{-s^2+2sx}}{s^{n+1}}\,ds.
\]

Interchange sum and integral:
\[
= \frac{1}{2\pi i}\oint\frac{e^{-s^2+2sx}}{s}\sum_{n=0}^{\infty}\left(\frac{t}{s}\right)^n\,ds = \frac{1}{2\pi i}\oint\frac{e^{-s^2+2sx}}{s-t}\,ds,
\]
where the geometric series converges for $|t/s| < 1$ (choose $C$ large enough).

By Cauchy's integral formula (the pole at $s = t$ is inside $C$):
\[
\frac{1}{2\pi i}\oint\frac{e^{-s^2+2sx}}{s-t}\,ds = e^{-t^2+2xt}.
\]

Therefore:
\[
\sum_{n=0}^{\infty}\frac{H_n(x)}{n!}t^n = e^{-t^2+2xt}. \quad
\]

\end{document}
